% ch4_hypotheses.tex — 第四章 研究假设
\section{研究假设}

本研究提出四个核心假设（H1--H4），形成"预测效度$\to$调节效应$\to$行为意向预测$\to$增量效度"的递进检验体系。H1/H2以全样本（$N=1{,}157$，ATT和Rating均完整）检验ATT对Rating的预测及其边界条件；H3/H4以明确表达行为意向的子样本（$N=298$，占总样本22.4\%）检验ATT对BI的预测力及其超越Rating的增量效度。消费者多属性态度通过以下公式计算：

\begin{equation}
\text{ATT} = \sum_{i}(b_i \times e_i)
\end{equation}

其中$b_i$为消费者对第$i$个产品属性信念的确信程度（0.4--1.0连续值），$e_i$为该信念的情感极性（$+1$积极/$0$中性/$-1$消极），ATT为所有信念的加权求和。

\subsection{ATT预测效度与调节效应假设H1和H2}

\subsubsection{H1：多属性态度正向预测评论评分}

根据Fishbein多属性态度理论\cite{fishbein1975}，消费者对产品的总体态度是其对各属性信念的加权综合，态度越积极，消费者对产品的整体评价应越正面。评论评分（Rating）是消费者在完成购买并使用产品后对其整体满意度的客观记录，是其产品态度最直接的外显指标。与来自评论文本的LLM提取数据不同，Rating来自平台结构化字段，消费者在提交评分时不受文本表达能力的限制，能够更直接地反映其整体态度倾向。

若LLM从评论文本中提取的多属性态度（ATT）能够有效捕捉消费者的真实态度，则ATT应能正向预测消费者的客观评分行为。这一关系的成立，既验证了Fishbein多属性态度公式的有效性，也证明了LLM提取方法的预测效度（predictive validity）——即测量工具能够预测与构念理论相关的外在指标。

需要指出，Rating数据在本研究样本中呈典型的J型分布（5星占比72.2\%，1星占6.3\%），这一偏态特征意味着ATT对评分高分段内部变异的解释力可能受到天花板效应的制约，但对区分低评分与高评分评论仍应保持较强的区分能力。鉴于此，H1的分析阶段将采用有序Logistic回归处理评分的有序离散特征\cite{duan2023,kim2022}，以OLS线性回归作为稳健性检验。

\textbf{H1：LLM提取的消费者多属性态度（ATT）正向预测评论评分（Rating）。}

\subsubsection{H2：评论者经验对ATT-Rating关系的调节效应}

在确认ATT能够正向预测Rating（H1）的基础上，一个进一步的问题是：这种预测关系的强度是否因评论者而异？Rocklage et al.\cite{rocklage2023}的研究指出，个体通过直接体验（direct experience）而非间接信息形成的态度，具有更高的可及性（accessibility）和清晰度，因而与外显行为的一致性显著更强。对于在线评论而言，评论者的历史经验越丰富，其对产品属性的判断越精准，形成的多属性态度越清晰稳定。

Mitchell et al.\cite{mitchell2022}的多领域实证分析进一步将经验与认知能力联系起来：经验丰富的消费者拥有更精细的产品知识图式（schema），能够更准确地识别和评估产品属性信念，其态度判断与外显行为（如评分）之间的对应关系更为紧密。值得注意的是，消费者专业性（consumer expertise）与消费者知识（consumer knowledge）存在细微差别：前者强调通过实践积累的主观能力和经验，后者偏重客观认知存量；本研究以评论数量代理的是专业性中的"品类熟悉度"维度，而非客观知识量。

据此，评论者经验越丰富，其多属性态度与评分行为的一致性应越强，即评论者经验正向调节ATT对Rating的预测关系。

然而，值得注意的是，高经验评论者可能对产品形成更差异化、更严苛的评价标准，从而对同一产品在态度上给予更高要求，但在评分上相对保守（即"天花板效应"）；或其多属性态度已高度内化，文本表达与评分之间的对应关系反而趋于稳定但幅度更小。因此，评论者经验对ATT--Rating关系的调节方向仍是一个有待检验的实证问题，本研究将如实报告交互项的方向与显著性。

\textbf{H2：评论者经验（Reviewer Experience）调节ATT对Rating的预测关系（探索性假设，交互方向有待检验）。}

\subsection{ATT对行为意向的预测力与增量效度假设H3和H4}

\subsubsection{H3：多属性态度正向预测行为意向}

计划行为理论\cite{ajzen1991}是行为科学领域最具影响力的理论框架之一，该理论将态度（attitude toward the behavior）定义为行为意向最直接的前因变量之一，认为个体对某一行为持有的态度越积极，其执行该行为的意向越强。Sheeran \& Webb\cite{sheeran2022}对十年间研究的元分析表明，态度与行为意向之间存在中等强度的正向关系（$r \approx 0.50$），这一关系已在跨文化、跨行为类型的大量研究中得到验证。

在在线评论情境中，消费者完成购买并使用产品后，其基于实际使用体验形成的产品态度（ATT）应正向影响其未来行为意向（BI），包括再次购买意向和推荐他人意向。这一关系在消费者行为研究中已有充分的理论依据\cite{fishbein1975}。本研究将上述经典理论应用于评论文本情境，检验从评论文本中LLM提取的ATT是否能够正向预测从同一文本提取的BI。

\textbf{H3：消费者多属性态度（ATT）正向预测行为意向（BI）。}

\textit{注：H3和H4的检验样本为明确表达行为意向的评论者子样本（$N=298$，占总样本22.4\%），ATT与BI均从同一评论文本提取，存在同源方差（CMV）问题；该子样本存在选择偏差（BI存在组ATT均值显著更高，$t=4.60$，$p<0.001$），研究结论需审慎解读，详见局限性章节。关于替代BI测量方案的探索与排除，详见8.5.1节。}

\subsubsection{H4：多属性态度提供超越评分的增量预测力}

多属性态度理论的核心价值在于通过对各具体属性信念的分解，捕捉消费者态度的复杂结构。若ATT仅是Rating的代理变量（即两者测量的是完全相同的信息），则在控制Rating后，ATT对BI的预测力应消失。然而，Fishbein多属性态度理论\cite{fishbein1975}认为，多属性分解提供了单一总体评价之外的额外信息：消费者可能在评分上给出5星，但其多属性态度中同时包含若干负面信念（如"包装简陋但效果好"），这些细粒度信息能够更准确地预测其未来是否会再次购买或推荐。

若多属性分解确实提供了额外信息，则即使在控制Rating（代表消费者的总体态度水平）之后，ATT仍应对BI具有显著的增量预测力（incremental predictive validity）。这一假设的检验采用层次回归分析，比较加入ATT前后的$R^2$变化量（$\Delta R^2$）。

\textbf{H4：在控制Rating后，ATT对BI仍具有显著的增量预测力（$\Delta R^2>0$），证明多属性态度提供了超越单一评分的信息增量。}
